\subsubsection{Hilbert 空间的有界逆定理的证明}\label{subsec:hilbert-bounded-inverse-proof}

本节证明在 Section~\ref{sec:open-mapping-theorem} 中承诺的结果.

\begin{Theorem}
\label{thm:hilbert-bounded-inverse}
若 $X$ 和 $Y$ 为 Hilbert 空间, $A$ 为从 $X$ 到 $Y$ 的连续双射, 则 $A$ 的逆有界.
\end{Theorem}

\begingroup
\postdisplaypenalty=10000
\renewcommand{\endProof}{\par\nobreak\hfill$\blacksquare$\par\vspace{0.1cm}}
\begin{Proof}
由于 $A=A^{**}$, Problem~\ref{ex:hilbert-adjoint-bounded-inverse-equivalence} 表明, 只需证明 $A^*$ 有有界逆. 由于 $A$ 的零空间是平凡的, $A^*$ 的值域在 $X$ 中稠密. 因而只需证明存在 $\delta>0$, 使得对于所有 $y\in Y$ 均有
\begin{equation}\label{eq:hilbert-bounded-inverse-adjoint-lower-bound}
\|A^*y\|\geq\delta\|y\|
\end{equation}
反设不然, 则存在序列 $y^n$, 使得
\begin{equation}\label{eq:hilbert-bounded-inverse-adjoint-normalized-sequence}
\|A^*y^n\|=1
\end{equation}
且
\begin{equation}\label{eq:hilbert-bounded-inverse-unbounded-sequence}
\|y^n\|\to\infty.
\end{equation}
但现在, 对任意 $f\in Y$, 利用 $A$ 为满射这一事实, 令 $x$ 为 $Ax=f$ 的解. 那么
\begin{equation}\label{eq:hilbert-bounded-inverse-weak-bound}
|(y^n,f)|=|(y^n,Ax)|=|(A^*y^n,x)|\leq\|x\|;
\end{equation}
即序列 $y^n$ 弱有界. 由一致有界原理 (Theorem~\ref{thm:uniform-boundedness}), $y^n$ 的范数必有界, 产生矛盾.
\end{Proof}
\endgroup
